\subsection{Characters and Grothendieck Groups}

Denote by \(\theta_{G}\) the K-algebra homomorphism from \(K\otimes_{\mathbf{Z}}\mathrm{R}_{\mathrm{K}}(\mathrm{G})\) to \(\mathcal{Z}_{\mathrm{K}}(\mathrm{G})\) that sends \(1\otimes[\pi]\) to \(\chi_{\pi}\) for every finite-dimensional representation \(\pi\) .

PROPOSITION 8. --- a) The homomorphism \(\theta_{\mathrm{G}}\) is an isomorphism from \(\mathrm{K} \otimes_{\mathbf{Z}} \mathrm{R}_{\mathrm{K}}(\mathrm{G})\) to \(\mathcal{Z}_{\mathrm{K}}(\mathrm{G})\) .

\begin{enumerate}
\def\labelenumi{\alph{enumi})}
\setcounter{enumi}{1}
\tightlist
\item
  Suppose that K has characteristic 0. Then \(\theta_{\mathrm{G}}\) defines an isomorphism from \(\mathrm{R}_{\mathrm{K}}(\mathrm{G})\) to the subring of \(\mathcal{Z}_{\mathrm{K}}(\mathrm{G})\) consisting of the linear combinations of the characters \(\chi_{\lambda}\) for \(\lambda\) in \(\widehat{\mathrm{G}}\) with integer coefficients.
\end{enumerate}

The family \(([\lambda])_{\lambda \in \widehat{\mathbf{G}}}\) is a basis of the \(\mathbf{Z}\) -module \(\mathrm{R}_{\mathrm{K}}(\mathrm{G})\) , and the family \((\chi_{\lambda})_{\lambda \in \widehat{\mathbf{G}}}\) is a basis of the K-vector space \(\mathcal{Z}_{\mathrm{K}}(\mathrm{G})\) . Assertions a) and b) follow (VIII, p.~411, Proposition 5).

